\chapter{Various}

\section{Intervals}
	\kactlimport{IntervalContainer.h}
	\kactlimport{IntervalCover.h}
	\kactlimport{ConstantIntervals.h}

\section{Misc. algorithms}
	\kactlimport{TernarySearch.h}
	\kactlimport{LIS.h}
	\kactlimport{FastKnapsack.h}

\section{Dynamic programming}
	\subsection{Quadrangle Inequality}
		Quadrangle Inequality: some cost function $C$ satisfies for all $l < r$: $C(l,r+1)-C(l, r) \ge  C(l+1,r+1) - C(l+1,r)$.
		Alternatively $C(a,c) + C(b,d) \le C(a,d) + C(b,c)$ for all $a \le b \le c \le d$. Still holds for $C'(l,r)=C(l,r)+F(l)+G(r)$.
		Example: sub-matrix sum of $(l, l)$ to $(r, r)$. 
		Define $\text{opt}(r)=$(leftmost or rightmost)$l(< r)$ that minimize $C(l,r)$, then $\text{opt}$ is non-decreasing.
		$\text{opt}$ is also the decision point of DP transitions. Consider writing assertions of monotonicity of $\text{opt}$ 
		if inequality is difficult to prove.
		This gives a sufficient condition for applying Divide and Conquer DP, Knuth DP or 1D1D DP.
	\kactlimport{KnuthDP.h}
	\kactlimport{DivideAndConquerDP.h}
	\kactlimport{1D1D.h}
	\subsection{Sum over Subset}
		Keep set $S(\text{mask}, i)$ as sum of (subset of first $i$ bits of $\text{mask}$ + remaining bits of $\text{mask}$).
		Loop over $i$ first, then $\text{mask}$. If $\text{mask}$'s $i^{th}$ bit is $1$, 
		additionally include $S(\text{mask} \oplus 2^i, i-1)$.
		Get rid of the second dimension for memory optimization.
	% TODO: slope trick
	% TODO: knapsack, subset sum and (max, +) convolution
	% TODO: Alien's trick
	% TODO: Dynamic Submask Count dp
	% TODO: O(nk) subtree dp
\vspace{10px}

\section{Debugging tricks}
	\begin{itemize}
		\item \verb@signal(SIGSEGV, [](int) { _Exit(0); });@ converts segfaults into Wrong Answers.
			Similarly one can catch SIGABRT (assertion failures) and SIGFPE (zero divisions).
			\verb@_GLIBCXX_DEBUG@ failures generate SIGABRT (or SIGSEGV on gcc 5.4.0 apparently).
	\end{itemize}
	
\vspace{10px}

\section{Optimization tricks}
	\verb@__builtin_ia32_ldmxcsr(40896);@ disables denormals (which make floats 20x slower near their minimum value). Needs 64-bit OS.
	\subsection{Bit hacks}
		\begin{itemize}
			\item \verb@x & -x@ is the least bit in \texttt{x}.
			\item \verb@for (int x = m; x; ) { --x &= m; ... }@ loops over all subset masks of \texttt{m} (except \texttt{m} itself).
			\item \verb@c = x&-x, r = x+c; (((r^x) >> 2)/c) | r@ is the next number after \texttt{x} with the same number of bits set.
			\item \verb@rep(b,0,K) rep(i,0,(1 << K))@ \\ \verb@  if (i & 1 << b) D[i] += D[i^(1 << b)];@ computes all sums of subsets.
		\end{itemize}
	\subsection{Pragmas}
		\begin{itemize}
			\item \lstinline{#pragma GCC optimize ("Ofast")} will make GCC auto-vectorize loops and optimizes floating points better.
			\item \lstinline{#pragma GCC target ("avx2")} can double performance of vectorized code, but causes crashes on old machines. Consider putting loops over arrays in a separate method.
			\item \lstinline{#pragma GCC optimize ("trapv")} kills the program on integer overflows (but is really slow).
		\end{itemize}
	\kactlimport{FastMod.h}
	\kactlimport{FastInput.h}
	\kactlimport{BumpAllocator.h}
	\kactlimport{SmallPtr.h}
	\kactlimport{BumpAllocatorSTL.h}
	% \kactlimport{Unrolling.h}
	\kactlimport{SIMD.h}
